\chapter{The rest of rcore: saving, the window, other input}

rcore is a third of the library. Chapters 2 to 5 covered the parts you need on
day one; this is the rest, and it is mostly the plumbing a finished game needs.

\section{Saving and loading}

The simplest save file that works is text:

\begin{lstlisting}
// writing
const char *text = TextFormat("score=%i\nlevel=%i\n", score, level);
SaveFileText("savegame.txt", (char *)text);

// reading
if (FileExists("savegame.txt"))
{
    char *text = LoadFileText("savegame.txt");   // heap allocated
    // ... parse it ...
    UnloadFileText(text);                        // you must free it
}
\end{lstlisting}

For binary data --- a struct, a level, a screenshot --- the pair is
\lstinline|SaveFileData| / \lstinline|LoadFileData|, and the loader hands you the
size:

\begin{lstlisting}
int size = 0;
unsigned char *data = LoadFileData("level.bin", &size);
UnloadFileData(data);
\end{lstlisting}

\gotcha{Every \lstinline|Load...| in this family allocates. Every one has a
matching \lstinline|Unload...|. This is the same discipline as textures, and it
is easier to forget here because nothing appears on screen to remind you.}

\section{Paths, and the one that actually matters}

\begin{center}
\begin{tabular}{@{}ll@{}}
\toprule
\lstinline|GetWorkingDirectory()| & where the program was \emph{launched} from \\
\lstinline|GetApplicationDirectory()| & where the executable \emph{sits} \\
\lstinline|ChangeDirectory(path)| & move the working directory \\
\lstinline|GetFileName(path)| & \lstinline|"boss.png"| \\
\lstinline|GetFileNameWithoutExt(path)| & \lstinline|"boss"| \\
\lstinline|GetFileExtension(path)| & \lstinline|".png"| \\
\lstinline|GetDirectoryPath(path)| & the folder part \\
\lstinline|IsFileExtension(path, ".png")| & \\
\lstinline|FileExists / DirectoryExists| & \\
\lstinline|GetFileLength(path)| & size in bytes \\
\lstinline|LoadDirectoryFiles(path)| & list a folder \\
\bottomrule
\end{tabular}
\end{center}

Most of those only parse the string and never touch the disk, so they are safe to
call freely.

The distinction between the two directory functions is the source of the classic
``it works in my IDE but not when I double-click it'' bug. Relative paths resolve
against the \emph{working} directory, which is not necessarily where your game
lives. The fix is one line at startup:

\begin{lstlisting}
ChangeDirectory(GetApplicationDirectory());
\end{lstlisting}

or this project's helper, \lstinline|SearchAndSetResourceDir("resources")|.

\section{Dropped files}

\begin{lstlisting}
if (IsFileDropped())
{
    FilePathList dropped = LoadDroppedFiles();

    for (unsigned int i = 0; i < dropped.count; i++)
        TraceLog(LOG_INFO, "dropped: %s", dropped.paths[i]);

    UnloadDroppedFiles(dropped);      // always, or you leak
}
\end{lstlisting}

Twenty lines and your level editor can open files by drag and drop. It is the
cheapest tooling win in the library.

\section{Random numbers}

\begin{lstlisting}
SetRandomSeed(1234);                       // makes a run repeatable
int roll = GetRandomValue(1, 6);           // inclusive at both ends

int *deck = LoadRandomSequence(52, 1, 52); // shuffled, NO repeats
UnloadRandomSequence(deck);
\end{lstlisting}

\lstinline|LoadRandomSequence| is the one people reimplement badly. Any time you
need ``pick these in a random order without repeating'' --- cards, spawn points,
a shuffled playlist --- it is already there.

Seeding deliberately is worth the habit: a fixed seed makes a bug reproducible,
and it is how ``daily challenge'' levels work.

\section{Logging}

\begin{lstlisting}
TraceLog(LOG_INFO,    "level %i loaded", level);
TraceLog(LOG_WARNING, "texture missing, using placeholder");
TraceLog(LOG_ERROR,   "save failed");

SetTraceLogLevel(LOG_WARNING);    // quieten raylib's own chatter
\end{lstlisting}

raylib logs a lot at \lstinline|LOG_INFO|, which is genuinely useful while
learning --- it is how you confirm a shader compiled or a texture loaded. Turn it
down to \lstinline|LOG_WARNING| for a release.

\section{Managing the window at run time}

\begin{center}
\begin{tabular}{@{}ll@{}}
\toprule
\lstinline|ToggleFullscreen()| & a real mode switch; can be slow \\
\lstinline|ToggleBorderlessWindowed()| & fake fullscreen; instant, usually preferred \\
\lstinline|MaximizeWindow() / RestoreWindow()| & needs \lstinline|FLAG_WINDOW_RESIZABLE| \\
\lstinline|SetWindowSize(w, h) / SetWindowPosition(x, y)| & \\
\lstinline|SetWindowTitle(text)| & \\
\lstinline|SetWindowIcon(image)| & \\
\lstinline|IsWindowResized()| & true for the one frame it changed \\
\lstinline|IsWindowFocused()| & pause the game when it goes false \\
\bottomrule
\end{tabular}
\end{center}

And the monitor functions, for an options menu that offers real resolutions:

\begin{lstlisting}
int m = GetCurrentMonitor();
GetMonitorName(m);  GetMonitorWidth(m);  GetMonitorHeight(m);  GetMonitorRefreshRate(m);
\end{lstlisting}

\gotcha{\lstinline|GetScreenWidth()| is the window in logical units;
\lstinline|GetRenderWidth()| is the actual framebuffer in pixels. On a HighDPI
display they differ by a factor of two, and mixing them puts your HUD in the
wrong place on exactly the machines you do not own.}

\section{Gamepads}

\begin{lstlisting}
if (IsGamepadAvailable(0))
{
    float x = GetGamepadAxisMovement(0, GAMEPAD_AXIS_LEFT_X);   // -1 .. +1
    if (IsGamepadButtonPressed(0, GAMEPAD_BUTTON_RIGHT_FACE_DOWN)) Jump();
}
\end{lstlisting}

The button names are described by position, not by letter, because A on an Xbox
pad is B on a Nintendo one: \lstinline|RIGHT_FACE_DOWN| is the bottom button of
the right cluster whatever the pad calls it.

\gotcha{Sticks never read exactly zero when released. Always apply a dead zone:

\lstinline|if (fabsf(x) < 0.15f) x = 0.0f;|

Without it your character drifts slowly across the room while nobody is touching
anything.}

\section{Touch and gestures}

\begin{lstlisting}
int count = GetTouchPointCount();
Vector2 p = GetTouchPosition(0);

int g = GetGestureDetected();     // GESTURE_TAP, GESTURE_SWIPE_LEFT, GESTURE_PINCH_OUT...
\end{lstlisting}

On a desktop the mouse counts as one touch point, so tap, hold, drag and swipe
all work without a touchscreen --- which means you can develop and test the mobile
controls of a game on your laptop.

\section{Clipboard and compression}

\begin{lstlisting}
SetClipboardText("share this");
const char *text = GetClipboardText();

unsigned char *packed = CompressData(raw, rawSize, &packedSize);
unsigned char *back   = DecompressData(packed, packedSize, &rawSize);
char *b64 = EncodeDataBase64(raw, rawSize, &outSize);
\end{lstlisting}

Compression plus base64 is how you make a save file or a level that a player can
copy and paste to a friend as a block of text.

\gotcha{These return raylib-allocated memory. Free it with
\lstinline|MemFree()|, not \lstinline|free()| --- raylib can be built with a
custom allocator, and mixing them crashes only on some builds.}

\tryit{Module 4 writes and reads a real save file, lists the path helpers with
their output, shuffles a deck, compresses a buffer and copies to your clipboard.
Module 1 covers the window and monitor side, and module 2 the input
devices.}{./cheatsheet/build.sh 4 \quad 1 \quad 2}
